% Curriculum vitae, public English version: the PDF the website serves at
% /downloads/konstantinos-papadogiannis-cv.pdf. Build with ./build.sh.
%
% The same content as the printed CV handed out in person, whose QR code opens
% the website's résumé page, so whoever scans it downloads what they hold.
% Unlike the printed one it carries only the contact details the website
% already shows (city, email, LinkedIn, GitHub), no QR code and no link to a
% project repository. See README.md.
\documentclass[english]{kpcv}

% Keep the PDF free of build details: no TeX banner, no source file name and
% no copy of the build information dictionary.
\pdfvariable suppressoptionalinfo \numexpr 1 + 2 + 4 + 8 \relax
% Plain objects instead of compressed object streams, so the link targets and
% the document information can be checked in the file itself.
\pdfvariable objcompresslevel 0

\hypersetup{
  pdftitle    = {Konstantinos Papadogiannis, Curriculum Vitae},
  pdfauthor   = {Konstantinos Papadogiannis},
  pdfsubject  = {Curriculum Vitae},
  pdfkeywords = {Konstantinos Papadogiannis, CV, junior software engineer, backend, full-stack,
                 Go, Java, Spring Boot, Python, TypeScript, React, Kubernetes, Docker, DevOps, cloud,
                 Hyperledger Fabric, machine learning, NLP, Harokopio University},
  pdflang     = {en-GB},
}

% Darker than the class defaults, as on the printed version
\definecolor{kpmuted}{HTML}{444B55}
\definecolor{kprule}{HTML}{A3ABB6}

% One contact item in the header: icon, then the detail
\newcommand*{\cvcontact}[2]{\mbox{\cvicon{#1}\enspace#2}}
\newcommand*{\cvgap}{\hspace{1.3em}}

\name{Konstantinos}{Papadogiannis}
\headline{Junior Software Engineer\sep Backend\sep Cloud \& DevOps\sep Machine Learning}
\contact{%
  \cvcontact{map-marker-alt}{Athens, Greece}\cvgap
  \cvcontact{envelope}{\href{mailto:konstantinospapadogiannis.dev@gmail.com}{konstantinospapadogiannis.dev@gmail.com}}\\
  \cvcontact{linkedin-in}{\href{https://www.linkedin.com/in/konstantinos-papadogiannis/}{in/konstantinos-papadogiannis}}\cvgap
  \cvcontact{github}{\href{https://github.com/KonstantinosPapadogiannis}{github.com/KonstantinosPapadogiannis}}}

% Header aligned left, as on the printed version, which has its QR code on the right
\makeatletter
\renewcommand*{\makecvheader}{%
  \begingroup
  \noindent\raggedright
  {\fontsize{24}{27}\selectfont
    {\kp@boldfont\kp@firstname\ \kp@lastname}\par}%
  \vspace{4pt}%
  {\fontsize{9.4}{12.5}\selectfont\bfseries\color{kpaccent}%
    \addfontfeatures{LetterSpace=5}\kp@headline\par}%
  \vspace{6pt}%
  {\fontsize{9.4}{14}\selectfont\color{kpmuted}\kp@contact\par}%
  \endgroup
  \vspace{2pt}}
\makeatother

% A little more air between sections, entries and bullets than kpcv.cls gives, as in cv-en-census.tex
\makeatletter
\renewcommand*{\cvsection}[1]{%
  \par\addvspace{10pt}%
  {\fontsize{10.2}{12}\selectfont\bfseries\color{kpaccent}%
    \addfontfeatures{LetterSpace=5}#1\par}%
  \nopagebreak\nointerlineskip\vspace{2.5pt}%
  {\color{kprule}\hrule height 0.5pt}%
  \nopagebreak\nointerlineskip\vspace{5pt}}
\makeatother
\setlist*[cvitems]{itemsep=1.6pt, topsep=2pt}
\let\kpcventry\cventry
\renewcommand*{\cventry}{\par\addvspace{7pt}\kpcventry}

% Skills as one run-in line each ("Programming: Go, ..."), so text extraction keeps label and values together
\renewcommand*{\cvskill}[2]{%
  \par\noindent\hangindent=1.2em\hangafter=1
  {\bfseries #1:}\enspace#2\par\addvspace{1.2pt}}

\begin{document}
\makecvheader

\cvsummary{Informatics and Telematics graduate (B.Sc.\ 2026) who builds backend services, web
applications and their cloud deployment in Go, Java/Spring Boot, Python, React/TypeScript and
Kubernetes. Thesis: a Hyperledger Fabric system whose browser wallet signs on the user's own device;
also machine learning and NLP projects measured against baselines.}

\cvsection{EDUCATION}
\cventry{B.Sc.\ in Informatics and Telematics}{2022 -- 2026}
  {Harokopio University of Athens}{Average 7.51 / 10}
\begin{cvitems}
  \item Bachelor's Thesis: 10 / 10
  \item Relevant coursework: DevOps Concepts and Tools~(10/10), Computer Networks~(9/10),
        Distributed Systems~(8/10)
\end{cvitems}

\cvsection{BACHELOR'S THESIS}
% Translation of the official Greek title (the thesis has no official English title)
\cventry{Development of a Decentralised Application (dApp) and Smart Contracts for the Digital
  Submission and Management of Subsidies}{Apr 2026 -- Sep 2026}
  {Written alongside a companion thesis on the network infrastructure; supervisor:
  Asst.~Prof.~Anargyros Tsadimas}{}
\begin{cvitems}
  \item Built the application layer of a seven-organisation Hyperledger Fabric network for EU farm
        subsidies: Go chaincode, backend APIs and React/TypeScript portals
  \item Developed a Manifest V3 browser wallet that generates each farmer's ECDSA P-256 key with the
        Web Crypto API, encrypts it with AES-GCM under a PBKDF2-derived key and signs transactions
        locally, with proposal checks that block blind signing
  \item Enforced a versioned subsidy policy with role-based access control (Keycloak OIDC) and private
        data collections
  \item Tested roles, state transitions and end-to-end flows, and benchmarked wallet signing on 2011
        and 2022 hardware
\end{cvitems}

\cvsection{SELECTED TECHNICAL PROJECTS}
\cventry{TradingAgent: ML decision-support system for crypto markets}{Oct 2026 -- Present}
  {Independent research project, team of two}{}
\begin{cvitems}
  \item Co-design a system in which a custom ML model issues calibrated probabilistic forecasts and a
        separate decision agent turns them into recommendations within cost and risk limits
  \item Evaluate gradient boosting, LSTM and transformer models with walk-forward backtesting on
        point-in-time data, pre-registered hypotheses and a locked hold-out period
\end{cvitems}

\cventry{OpenStack private cloud on Google Cloud}{Jun 2026}
  {Cloud Computing Management, individual}{}
\begin{cvitems}
  \item Installed DevStack on a GCP VM and hosted a three-tier app, working around GCP's
        floating-IP and routing restrictions
\end{cvitems}

\cventry{Social Media Monitoring Agent: search, NLP and summaries over tweets}{Apr 2026}
  {Natural Language Processing and Information Retrieval, team of two}{}
\begin{cvitems}
  \item Built hybrid BM25 and sentence-transformer search over 10,896 tweets in FAISS, evaluated with
        Precision@k and nDCG@k
  \item Classified emotions with a pre-trained Twitter-RoBERTa model (micro F1 0.71 against 0.63 for
        a TF-IDF baseline) and served RAG summaries with DistilBART through FastAPI
\end{cvitems}

\cventry{ErgoHub: freelance marketplace on Azure Kubernetes}{Feb 2025 -- Jul 2025}
  {DevOps Concepts and Tools, team of two}{}
\begin{cvitems}
  \item Built Jenkins pipelines and Ansible playbooks to test, containerise and deploy the app
  \item Deployed Spring Boot and React/Redux to AKS behind NGINX Ingress with TLS, PostgreSQL
        and MinIO
\end{cvitems}

\cvsection{SKILLS}
\cvskill{Programming}{Go, Java, Python, TypeScript, JavaScript, SQL, C, CUDA, Bash}
\cvskill{Backend}{Spring Boot, REST, gRPC, FastAPI, Hyperledger Fabric (chaincode, private data,
  Gateway), Keycloak (OIDC), JWT}
\cvskill{Frontend}{React, Redux, Tailwind CSS, Vite, browser extensions (Manifest V3), Web Crypto API}
\cvskill{Cloud and DevOps}{Docker, Kubernetes (AKS), Jenkins, Ansible, OpenStack, Google Cloud, Azure,
  NGINX, Linux, Git}
\cvskill{Data and ML}{PostgreSQL, PyTorch, Hugging Face Transformers, scikit-learn, pandas, FAISS}
\cvskill{Testing}{Go testing, Vitest, Playwright, JUnit, Jest}
\cvskill{AI tools}{Claude Code in daily development, with a custom code-review subagent}
\cvskill{Languages}{Greek (native), English (professional working proficiency)}

\end{document}
